% Separate analytical review increment. No automatic merge.
% Unique label prefix: inc:LP:v1:
\documentclass[11pt]{article}
\usepackage[T1]{fontenc}
\usepackage[utf8]{inputenc}
\usepackage{lmodern,amsmath,amssymb,amsthm,mathtools,microtype}
\usepackage[margin=0.95in]{geometry}
\usepackage{booktabs,array,xurl}
\usepackage[hidelinks]{hyperref}
\usepackage{fancyhdr}
\pagestyle{fancy}\fancyhf{}
\fancyhead[L]{\small Linear contrast and limiting persistence}
\fancyhead[R]{\small Review increment v1}
\fancyfoot[C]{\thepage}
\setlength{\headheight}{14pt}
\setlength{\emergencystretch}{3em}
\allowdisplaybreaks[2]
\numberwithin{equation}{section}
\newtheorem{theorem}{Theorem}[section]
\newtheorem{lemma}[theorem]{Lemma}
\newtheorem{proposition}[theorem]{Proposition}
\newtheorem{corollary}[theorem]{Corollary}
\theoremstyle{remark}\newtheorem{remark}[theorem]{Remark}
\newcommand{\E}{\mathbb E}
\newcommand{\R}{\mathbb R}
\newcommand{\dd}{\,d}
\newcommand{\norm}[1]{\left\lVert#1\right\rVert}
\DeclareMathOperator{\diag}{diag}
\title{Exact linear contrast and eventual permanent loss of probe benefit\
\large A separate analytical increment, not initialized Theorem A}
\author{Research proof material for independent review}
\date{Version 1 --- not merged into the checkpoint}
\begin{document}\maketitle

\begin{abstract}
We give two analytical statements with different scopes. First, a commuting,
aligned two-layer linear initialization has an explicit diagonal logistic
solution at finite width or in the continuum. With diagonal cluster second
moments, every training cluster helps every probe. This extends and restates
P10 of the supplied foundations; it is not a newly discovered architecture
control. Second, the selected coherent weak-growth orbit of the defined
small-noise system has $x(s)\to+\infty$, is eventually strictly increasing,
and has exactly zero scaled probe benefit for all sufficiently late times.
The basic persistence conclusion follows immediately from finite downward
variation and the existing arbitrary-level hitting theorem. A separate
comparison using $x+v$ proves eventual monotonicity without the unsupported
moving-barrier step in the proposed outline. A final conditional diagonal
argument explains what additional finite-noise convergence would imply
persistence on growing windows. No numerical trajectory, interval arithmetic,
or saved-array certificate is a premise.
\end{abstract}

\section{Scope, dependencies, and normalization}
\label{inc:LP:v1:scope}

The original initialized Theorems A and B remain open. This file does not
prove isotropic capture, nonlinear passage of the full population, control
of a non-small bridge, or any finite-noise trajectory transfer. It does not
infer infinite-time behavior from an experiment ending at time $100$.

The linear statement is proved directly below and extends \cite{F}, P10,
from scalar-isotropic initial covariance to an arbitrary positive diagonal
one. The limiting statements import the following from \cite{P},
\texttt{pa:branch} and \texttt{pa:limitreversal}: the selected orbit exists
globally forward, $v>0$, $v>y$, $v>x$, its coordinates grow at most linearly
forward, and its input coordinate reaches every finite level greater than
$a_\rho$. We restate the exact equations used. Uniformity over compact
parameter sets additionally uses the smooth selected-branch dependence
proved in \cite{W}, Lemma \texttt{inc:AN06:v2:branchsmooth}.

Time $\tau$ for the normalized original flow satisfies
$\tau=\mu_1^2t_{\mathrm{phys}}$. The centered limiting time is
$s=\tau-\tau_h$, with weak half-mass at $s=0$. Limits of physical rates
therefore use $D_p^{\mathrm{phys}}/(\mu_1^2q^2)$, not
$D_p^{\mathrm{phys}}/q^2$.

\section{The commuting linear architecture control}
\label{inc:LP:v1:linear}
Let $(\Omega,\kappa)$ be a label measure space. It can be a finite set with
counting measure, or a continuum probability space. For square-integrable
$u,w:\Omega\to\R^d$, put
\[
 M=\int_\Omega w_\alpha u_\alpha^\top\dd\kappa(\alpha),
 \qquad f(X)=MX,\qquad G_0=\int u_{\alpha,0}u_{\alpha,0}^\top\dd\kappa.
\]
Let the training distribution be $\sum_p\pi_pP_p$, where $\pi_p>0$ and
$\sum_p\pi_p=1$. In this section use \emph{weighted} second moments
\[
 \mathsf A_p=\pi_p\E_{P_p}[XX^\top],\qquad A=\sum_p\mathsf A_p.
\]
This differs from the raw-cluster convention in P10 by the included
mixture weights. The loss is $\mathcal L=\tfrac12\E\norm{MX-X}^2$.
The Euclidean or $L^2(\kappa)$ gradient equations are
\begin{equation}\label{inc:LP:v1:label-flow}
 w_\alpha'=Ru_\alpha,\qquad u_\alpha'=R^\top w_\alpha,
 \qquad R=(I-M)A.
\end{equation}

\begin{theorem}[Diagonal aligned linear flow]
\label{inc:LP:v1:linear-theorem}
Suppose
\[
 u_{\alpha,0}=w_{\alpha,0}\quad\kappa\text{-a.e.},\qquad
 G_0=\diag(c_1,\ldots,c_d),\quad 0<c_r<1,
 \qquad A=\diag(a_1,\ldots,a_d),\quad a_r>0.
\]
Then the unique solution is global and remains aligned:
\[
 u_\alpha(\tau)=w_\alpha(\tau)=T(\tau)u_{\alpha,0},\qquad
 T_{rr}(\tau)=\sqrt{\ell_r(\tau)/c_r},
\]
where
\begin{equation}\label{inc:LP:v1:logistic}
 \ell_r'=2a_r\ell_r(1-\ell_r),\qquad
 \ell_r(\tau)=\frac{1}{1+(c_r^{-1}-1)e^{-2a_r\tau}},
 \qquad M(\tau)=\diag(\ell_r(\tau)).
\end{equation}
For every nonzero probe $\xi$ and finite $\tau\ge0$,
\begin{equation}\label{inc:LP:v1:linear-error}
 E_\xi=\frac12\sum_r(1-\ell_r)^2\xi_r^2,
 \qquad E_\xi'=-2\sum_ra_r\ell_r(1-\ell_r)^2\xi_r^2<0,
 \qquad E_\xi=E_{-\xi}.
\end{equation}
If each $\mathsf A_p$ is also diagonal in this same basis, then its exact
cluster-sourced contribution is
\begin{equation}\label{inc:LP:v1:linear-source}
 D_p^\tau(\xi)
 =-2\sum_r(\mathsf A_p)_{rr}\ell_r(1-\ell_r)^2\xi_r^2\le0.
\end{equation}
It is strictly negative when $\xi$ has a nonzero coordinate on which
$(\mathsf A_p)_{rr}>0$.
\end{theorem}
\begin{proof}
The bilinear map $(u,w)\mapsto M$ is continuous from
$L^2(\kappa)\times L^2(\kappa)$ into the finite-dimensional matrix space,
by Cauchy--Schwarz. The right-hand side of
\eqref{inc:LP:v1:label-flow} is therefore locally Lipschitz on that Hilbert
space. Construct $\ell_r,T$ by \eqref{inc:LP:v1:logistic}. Their denominators
are positive at all nonnegative times and $0<\ell_r<1$. For the proposed
aligned solution,
\[
 M=TG_0T=\diag(\ell_r),\qquad
 R=\diag(a_r(1-\ell_r)),\qquad T'=RT.
\]
The last identity follows by differentiating
$T_{rr}=\sqrt{\ell_r/c_r}$. Both label equations are satisfied since
$R=R^\top$. Local uniqueness identifies this explicitly global candidate
with the initialized flow. Differentiating $E_\xi$ proves
\eqref{inc:LP:v1:linear-error}; at finite times every coefficient is
positive, so every nonzero probe has strict descent.

For the cluster source write $R_p=(I-M)\mathsf A_p$. Directly
from the two untied label velocities,
\begin{equation}\label{inc:LP:v1:source-matrix}
 M'_p=R_pP_U+P_WR_p,\quad
 P_U=\int uu^\top\dd\kappa,\quad P_W=\int ww^\top\dd\kappa.
\end{equation}
On the constructed trajectory $P_U=P_W=M$. When $\mathsf A_p$ is diagonal,
$M'_p=2(I-M)\mathsf A_pM$. Contracting its probe velocity with
$-(I-M)\xi$ gives \eqref{inc:LP:v1:linear-source}. The mixture weight has
already been included in $\mathsf A_p$.
\end{proof}

\begin{remark}[Necessary qualifications]
The total diagonal moment alone does not imply each cluster helps. For
example, at an aligned state with $M=\diag(1/5,4/5)$, let
\[
 \mathsf A_1=\begin{pmatrix}1/2&49/100\\49/100&1/2\end{pmatrix},\qquad
 \mathsf A_2=\begin{pmatrix}1/2&-49/100\\-49/100&1/2\end{pmatrix}.
\]
Both matrices are positive definite and $A=I$, but
\eqref{inc:LP:v1:source-matrix} gives
$D_1((1,-1)/\sqrt2)=433/5000>0$ at that state. The diagonal cluster
assumption must therefore accompany \eqref{inc:LP:v1:linear-source}.
Similarly, positive $a_r$ is needed for strict descent of every direction;
a zero data eigenvalue leaves that coordinate unchanged.
\end{remark}

\begin{corollary}[SIM continuum and exact finite-width scope]
\label{inc:LP:v1:linear-SIM}
For normalized two-concept SIM,
\[
 \mathsf A_1=\frac12\diag(1+q^2,q^2),\quad
 \mathsf A_2=\frac12\diag(q^2,\rho^2+q^2),\quad
 A=\diag(1/2+q^2,\rho^2/2+q^2).
\]
Aligned isotropic continuum initialization has $G_0=S_0I/2$, so for
$0<S_0<2$ all conclusions of Theorem~\ref{inc:LP:v1:linear-theorem} hold.
They also hold at any finite width whose initial Gram matrix is exactly
positive diagonal, on these population distributions. Independent random
finite-width isotropic initialization is not exactly diagonal in general.
\end{corollary}
\begin{proof}
The Gaussian cluster moments are the displayed diagonal matrices. Uniform
unit directions in the plane have second moment $I/2$ by antipodal and
coordinate-exchange symmetry, giving $G_0=S_0I/2$. Apply the theorem.
\end{proof}

\begin{remark}[Positive-quadrant initialization]
For a uniform initial angle on $[0,\pi/2]$ and total squared mass $S_0$,
\[
 G_0=S_0\begin{pmatrix}1/2&1/\pi\\1/\pi&1/2\end{pmatrix}.
\]
This follows by integrating $\cos^2\alpha$, $\sin^2\alpha$ and
$\sin\alpha\cos\alpha$ with density $2/\pi$. It exhibits nonzero minor
entries and generally leaves the commuting class when the data eigenvalues
differ. It does not, on its own, prove reversal for that initialization.
The tied predecessor and the untied network must still be distinguished.
For the same mass the isotropic linear initial map is $S_0X/2$, whereas the
ReLU initial map is $S_0X/4$; matching initial outputs rescales linear mass,
without changing the monotonicity conclusion.
\end{remark}

\section{Selected limiting orbit and its already-proved properties}
\label{inc:LP:v1:orbit}
Fix $0<\rho<1$, $\lambda=\rho^2$, and
$K(z)=\phi(z)+z\Phi(z)$, where $\phi,\Phi$ are the standard Gaussian density
and distribution function. On the strong-learned coherent surface the
selected weak-growth orbit solves
\begin{align}
 n(s)&=(1+e^{-\lambda s})^{-1},\nonumber\\
 x'&=\frac\lambda2[nv-x+(1-n)y]\Phi(\rho x),\nonumber\\
 y'&=\frac12[-y-nK(u)+\rho(1-n)K(\rho x)],\nonumber\\
 u'&=\frac12[(1-n)(\phi(u)-\lambda u)-(nu+y)\Phi(u)],\nonumber\\
 v'&=\frac12[\rho K(\rho x)-\lambda v].
 \label{inc:LP:v1:orbit-equations}
\end{align}
It approaches $(n,x,y,u,v)=(0,a,a,u_*,a/\lambda)$ as $s\to-\infty$, with
$a=\rho K(\rho a)$ and
$\phi(u_*)-a\Phi(u_*)-\lambda u_*=0$.
The selected branch and its global continuation are supplied by \cite{P}.
In particular, for $s\ge0$ there is a finite constant $C$ such that
\begin{equation}\label{inc:LP:v1:imports}
 v>0,\quad v>y,\quad v>x,\qquad
 \max(|x|,|y|,|u|,|v|)\le C(1+s).
\end{equation}
The same source proves that $x$ reaches every finite level $\zeta>a$.
These are analytical orbit properties, not numerical observations.

Define
\[
 d=v-x>0,\qquad \epsilon=(1-n)(v-y)>0,\qquad
 P=\Phi(\rho x),\qquad
 g_\rho(x)=K(\rho x)/\rho-x>0.
\]
Direct subtraction of \eqref{inc:LP:v1:orbit-equations} yields
\begin{align}
 x'&=\frac\lambda2(d-\epsilon)P,\label{inc:LP:v1:x-equation}\\
 d'&=\frac\lambda2[g_\rho(x)-(1+P)d+\epsilon P],
 \label{inc:LP:v1:d-equation}\\
 v'&=\frac\lambda2[g_\rho(x)-d].\label{inc:LP:v1:v-equation}
\end{align}
Since $1-n\le e^{-\lambda s}$ for $s\ge0$, the imported growth bound gives
\begin{equation}\label{inc:LP:v1:epsilon-tail}
 0<\epsilon(s)\le2C(1+s)e^{-\lambda s},\qquad
 \frac\lambda2\int_s^\infty\epsilon(r)\dd r
 \le C e^{-\lambda s}(1+s+\lambda^{-1})=:Q(s)\longrightarrow0.
\end{equation}

\section{Permanent loss of the leading probe benefit}
\label{inc:LP:v1:persistence}
\begin{lemma}[Unboundedness with finite downward variation]
\label{inc:LP:v1:downward}
Let $z$ be locally absolutely continuous on $[0,\infty)$. If
$\int_0^\infty(z')_-<\infty$ and $\limsup_{s\to\infty}z(s)=+\infty$,
then $z(s)\to+\infty$. More precisely, for $t\ge s$,
\[
 z(t)\ge z(s)-\int_s^\infty(z')_-(r)\dd r.
\]
\end{lemma}
\begin{proof}
Integrate $z'= (z')_+-(z')_-$ on $[s,t]$ and discard its positive part.
Given a real threshold $H$, choose a sufficiently late $s$ for which
$z(s)>H+1$ and the tail integral is less than one. Such $s$ exists by the
unbounded limsup and convergence of the tail integral. The displayed
inequality implies $z(t)>H$ for every $t\ge s$.
\end{proof}

\begin{theorem}[Eventual permanent inactivity in the selected limit]
\label{inc:LP:v1:permanence}
For the selected coherent orbit and every $0<\rho<1$,
\[
 x(s)\longrightarrow+\infty.
\]
For every finite $\zeta$, there is a finite $S_\zeta$ such that
$x(s)>\zeta$ and
\begin{equation}\label{inc:LP:v1:eventual-zero}
 \mathcal E_\zeta(s)
 =\frac12(\zeta-x(s))_+^2
 -(\zeta-x(s))_+(\zeta-y(s))=0
 \qquad(s\ge S_\zeta).
\end{equation}
The same $S_\zeta$ works for every offset in a bounded interval with upper
endpoint $\zeta$. For $\zeta>a$, this follows after a genuine earlier
reversal; all of the negative scaled probe error at any earlier time is
eventually returned. In particular the earlier minimum lies below
$-(\zeta-a)^2/2$, while the eventual scaled error is exactly zero.
This theorem does not assert that re-entry is impossible immediately after
the first gate hit.
\end{theorem}
\begin{proof}
Equation \eqref{inc:LP:v1:x-equation} and $d,P>0$ imply
$(x')_-\le\lambda\epsilon/2$. Hence
\[
 x(t)\ge x(s)-Q(s)\qquad(t\ge s\ge0).
\]
The imported arbitrary-level hitting theorem and boundedness on compact
time intervals imply $\limsup_{s\to\infty}x(s)=+\infty$.
Lemma~\ref{inc:LP:v1:downward} proves the asserted limit. Choose a time
past which $x$ exceeds the largest offset. Both positive parts in
\eqref{inc:LP:v1:eventual-zero} then vanish. The strict earlier minimum
and reversal are the imported selected-orbit theorem, \cite{P},
\texttt{pa:limitreversal}. No eventual monotonicity estimate is required for
this argument.
\end{proof}

\begin{remark}[What is zero]
$\mathcal E_\zeta$ is the coefficient in the small-noise expansion
$E_{\xi_q}-1/2=q^2\mathcal E_\zeta+o(q^2)$ on a justified fixed window.
Equation \eqref{inc:LP:v1:eventual-zero} is exact in the defined limiting
system, not the assertion that the original isotropic network has zero
probe output at finite $q$. Also, the leading system is continued as an
ODE after gate loss; the training equations do not depend on $\zeta$.
Its diverging scaled angles prevent an automatic interchange of $q\to0$
and $s\to\infty$.
\end{remark}

\section{The additional eventual-monotonicity statement}
\label{inc:LP:v1:monotonic}
The proposed argument that a negative derivative whenever
$d>g_\rho(x)+\epsilon$ automatically yields tracking of that moving graph
is incomplete: the graph itself moves. The following proof uses a scalar
comparison with $x+v$ instead.

\begin{proposition}[Slow growth and eventual positive input velocity]
\label{inc:LP:v1:monotonic-theorem}
On the selected orbit,
\[
 x(s)=O(\sqrt{\log(2+s)}),\qquad x'(s)>0
 \quad\text{for all sufficiently large }s.
\]
The constants and starting time may depend on $\rho$ and the selected
orbit. No asymptotic equivalent, numerical gate time, or no-re-entry claim
at the first hit is asserted.
\end{proposition}
\begin{proof}
\emph{Bound the gap.}
The downward-variation estimate gives a finite lower bound
$x\ge x_-=x(0)-Q(0)$. Let
$B_\epsilon=\sup_{s\ge0}\epsilon(s)<\infty$. Since $g_\rho$ is positive
and decreasing, \eqref{inc:LP:v1:d-equation} gives
\[
 d'\le\frac\lambda2[g_\rho(x_-)+B_\epsilon-d].
\]
Thus $0<d\le D$, where
$D=\max\{d(0),g_\rho(x_-)+B_\epsilon\}$.

\emph{A scalar upper-growth comparison.}
Put $W=x+v=2x+d$. From
\eqref{inc:LP:v1:x-equation} and \eqref{inc:LP:v1:v-equation},
\begin{equation}\label{inc:LP:v1:sum}
 W'=\frac\lambda2[g_\rho(x)-(1-P)d-\epsilon P]
 \le\frac\lambda2g_\rho(x).
\end{equation}
Theorem~\ref{inc:LP:v1:permanence} implies $W\to\infty$. Choose $s_0$
such that $z=W-D\ge1$ and $x\ge0$ for $s\ge s_0$.
Since $x=(W-d)/2\ge z/2$ and, for $r\ge0$,
$g_\rho(r)\le\phi(\rho r)/\rho$, we have
\[
 z'\le c_1e^{-\lambda z^2/8},\qquad
 c_1=\rho\phi(0)/2.
\]
Define $F(z)=\int_0^z e^{\lambda r^2/8}\dd r$. The ordinary chain rule
and the preceding inequality give
\[
 F(z(s))\le F(z(s_0))+c_1(s-s_0)=:\mathcal Q(s).
\]
For $z\ge1$, integration over $[z-1,z]$ gives
$F(z)\ge e^{\lambda(z-1)^2/8}$. Consequently
\begin{equation}\label{inc:LP:v1:sqrtlog}
 x(s)\le\frac{W(s)}2
 \le\frac{D+1}{2}+\sqrt{\frac2\lambda\log\mathcal Q(s)}.
\end{equation}
This establishes the claimed logarithmic upper growth; eventual positivity
of $x$ gives the stated big-$O$ bound.

\emph{A polynomial lower bound beats the transient.}
For $r\ge0$,
\[
 g_\rho(r)=\frac1\rho\int_{\rho r}^\infty(t-\rho r)\phi(t)\dd t
 \ge\frac1\rho\phi(\rho r+2),
\]
by restricting to $t\in[\rho r+1,\rho r+2]$.
Set $A_0=\rho(D+1)/2+2$. Equation \eqref{inc:LP:v1:sqrtlog} and
$(a+b)^2/2\le a^2+b^2$ yield
\begin{equation}\label{inc:LP:v1:gaplower}
 g_\rho(x(s))\ge\frac{\phi(0)}\rho e^{-A_0^2}\mathcal Q(s)^{-2}
 =:c_2\mathcal Q(s)^{-2}\qquad(s\ge s_0).
\end{equation}
The exact gap equation gives
$d'+\lambda d\ge\lambda g_\rho(x)/2$. Integrating only over the last
unit interval, for $s\ge s_0+1$, and using that $\mathcal Q$ increases,
\[
 d(s)\ge\frac{c_2}{2}(1-e^{-\lambda})\mathcal Q(s)^{-2}.
\]
The right side decays polynomially, whereas
$\epsilon(s)\le2C(1+s)e^{-\lambda s}$ decays exponentially. Their ratio
therefore tends to infinity. Eventually $d>\epsilon$ and
\eqref{inc:LP:v1:x-equation} gives $x'>0$.
\end{proof}

\section{Uniform offsets and a conditional finite-noise implication}
\label{inc:LP:v1:transfer}
\begin{proposition}[A common limiting clearing time on compact parameters]
\label{inc:LP:v1:uniform-clear}
Let $\mathcal K$ be a compact parameter interval in $(0,1)$ on which the
selected branches depend continuously on $\rho$ on fixed finite windows.
Suppose the imported forward growth estimate is uniform on $\mathcal K$.
For every finite $\zeta_+$ there is a common finite $S$ such that
\[
 x_\rho(s)>\zeta_+\qquad(\rho\in\mathcal K,\ s\ge S).
\]
For the interval used by AN06, the assumed continuity is supplied by its
branch-smoothness lemma; uniform linear growth follows from the checkpoint's
Dini inequality and bounded initial states on the compact parameter set.
\end{proposition}
\begin{proof}
Write $\lambda_- =\min_{\mathcal K}\rho^2>0$ and
$\lambda_+=\max_{\mathcal K}\rho^2$. The uniform bound
$\epsilon_\rho(s)\le2C(1+s)e^{-\lambda_-s}$ supplies a uniform upper
bound on $\lambda_\rho\int_s^\infty\epsilon_\rho/2$ tending to zero.
Choose $s_0$ so that this tail bound is below $1/2$ afterward.
For every $\rho_0$, choose $s_{\rho_0}\ge s_0$ with
$x_{\rho_0}(s_{\rho_0})>\zeta_++2$, possible by divergence.
Continuity gives a neighborhood of $\rho_0$ on which
$x_\rho(s_{\rho_0})>\zeta_++1$. The downward-variation bound then keeps
$x_\rho>\zeta_++1/2$ for every later time in that neighborhood. A finite
subcover of $\mathcal K$ and the maximum of its times provide $S$.
\end{proof}

\begin{proposition}[A conditional growing-window transfer]
\label{inc:LP:v1:diagonal-transfer}
Let $\mathcal I$ be a compact set of scaled offsets and let $S$ be a common
clearing time for a limiting family, including any stated compact $\rho$
range. Consider an \emph{actual initialized} positive-noise family with
$S_0=S_0(q)$, and an analytically specified centered time $t_h(q)$.
Put
\[
 F_q(s,\zeta)=
 \frac{E_{\xi_q}(t_h(q)+s/\mu_1^2)-1/2}{q^2},\qquad
 \xi_q=(\sin(q\zeta),-\cos(q\zeta)).
\]
Suppose, as an additional analytically proved hypothesis, that for every
finite $T>S$ the full original flow satisfies
\begin{equation}\label{inc:LP:v1:compact-transfer}
 \sup_{\zeta\in\mathcal I,\ S\le s\le T}
 |F_q(s,\zeta)-\mathcal E_\zeta(s)|\longrightarrow0
 \quad(q\downarrow0),
\end{equation}
with uniformity over any included $\rho$ range. Then there are
$T_q\to\infty$ and $\varepsilon_q\to0$ such that
\[
 \left|E_{\xi_q}(t_h(q)+s/\mu_1^2)-\frac12\right|
 \le\varepsilon_q q^2
 \quad(S\le s\le S+T_q,\ \zeta\in\mathcal I).
\]
In particular the proposed lower bound follows. No rate of growth of
$T_q$ is supplied by this proposition.
\end{proposition}
\begin{proof}
The limiting observable is zero for $s\ge S$. For each positive integer
$j$, use \eqref{inc:LP:v1:compact-transfer} with $T=S+j$ to choose
$q_j>0$ so that the supremum is at most $1/j$ whenever $0<q\le q_j$.
Choose the thresholds strictly decreasing to zero, with
$q_{j+1}\le q_j/2$. For $0<q\le q_1$, let
$j(q)=\max\{j:q\le q_j\}$. It is finite, tends to infinity as $q\to0$,
and its defining estimate gives the result with
$T_q=j(q)$ and $\varepsilon_q=1/j(q)$.
\end{proof}

\begin{remark}[What this conditional statement does not supply]
Theorem A's error/rate control only on its crossing windows is weaker than
\eqref{inc:LP:v1:compact-transfer}. Initialized selection, splitting,
background contributions, and finite-noise evolution must still establish
this additional compact-window convergence. The diagonal argument is an
analytical implication, not a substitute for those estimates or an explicit
long-time approximation. For a broad-bridge limiting population, one first
needs its own limiting persistence theorem.

The left endpoint must lie after an analytically justified limiting clearing
or near-recovery time, not just at an arbitrary fixed delay after the
minimum. If the limiting coefficient is still negative there, a bound
$E\ge1/2-o(q^2)$ cannot hold at that time.
\end{remark}

\begin{remark}[Every label need not clear]
To prove a lower bound on probe error, it is enough to control the
remaining benefit:
\[
 E_\xi=\frac12-\xi^\top f(\xi)+\frac12\norm{f(\xi)}^2
 \ge\frac12-[\xi^\top f(\xi)]_+.
\]
Thus $[\xi^\top f(\xi)]_+\le\varepsilon q^2$ suffices. Requiring every
original label to become probe-inactive is unnecessary and, for the exact
finite-time isotropic continuum, is incompatible with the angular
surjectivity deduction in \cite{R}. Positive mass may remain in every open
angular arc even while its observable contribution is negligible.
\end{remark}

\section{Status for Theorems A and B}
\label{inc:LP:v1:status-final}
The linear architecture control is already supported by the foundations;
this file states a wider diagonal-initialization version and records the
necessary per-cluster assumptions. The selected coherent limiting orbit
now has an analytical eventual-persistence conclusion in addition to its
reversal conclusion. These are not initialized theorems for the original
ReLU population.

An empirical statement of persistence through physical time $100$ should
not be renamed infinite-time permanence. A finite-small-parameter theorem
cannot be promoted to the canonical point without bounds that actually
include $q=0.05$ and $S_0=2\cdot10^{-4}$. Conversely, the non-small bridge
amplitude enhancement is not logically required for a first initialized
small-shape reversal theorem. It remains important if the theorem is
advertised as covering the canonical enhanced trajectory.

The exact decay rate of the original training loss, the numerical
$1.10$ time delay, finite-width fluctuations, and all empirical hypothesis
outcomes are outside the proofs in this increment.

\section*{Review and integration contract}
This file and its companion review note are new, separate review material.
No foundation, notebook, appendix, handoff, or theory checkpoint has been
modified. After independent checking and explicit approval, the linear
extension belongs with P10 (or its main-paper statement), and the limiting
persistence theorem belongs after \texttt{pa:limitreversal}. The growing-window
transfer remains an explicitly conditional corollary. No automatic merge is
requested.

\begin{thebibliography}{9}
\bibitem{F} \emph{ReLU Population Foundations}, version 2.
User-supplied \path{RELU_POPULATION_FOUNDATIONS_v2.tex}, P10, Section 16.
\bibitem{P} \emph{Analytical Progress toward Initialized Population Reversal},
checkpoint 1.0. User-supplied \path{THEOREM_A_ANALYTICAL_PROGRESS.tex},
\texttt{pa:branch}, \texttt{pa:atomicrates}, \texttt{pa:limitreversal}.
\bibitem{W} \emph{Learned Limiting-Orbit Witnesses}, version 2.
User-supplied \path{AN06_learned_orbit_witnesses_v2.tex}, smooth branch
dependence and common-window constructions.
\bibitem{R} \emph{Reviewer Addendum to AN02 and AN06}, version 1.
User-supplied \path{AN02_AN06_review_addendum_v1.tex}, finite-time angular
surjectivity and background-budget distinctions.
\end{thebibliography}
\end{document}
